\chapter{曲線のねじれ具合}

これまでいくつかの曲線を見てきました。
曲線の曲がり具合としての「曲率」をうまく定義できました。
では「ねじれ」というものはどんなものでしょうか？
見ていきましょう。



\section{曲線はどれだけ「平面的」から離れているか？}

例えば円はねじれた曲線でしょうか？直感的にそうではなさそうです。
放物線はねじれているでしょうか？これもそんな気はしません。
ところが常螺旋はねじれている気がします。

それぞれの単位接ベクトル、主法線ベクトル、従法線ベクトルを書きだすと、円の場合
\[
  \begin{cases}
    e_1(s)=(-\sin(s/r),\cos(s/r),0)\\
    e_2(s)=(-\cos(s/r),-\sin(s/r),0)\\
    e_3(s)=(0,0,1)
  \end{cases}
\]
常螺旋の場合
\[
  \begin{cases}
    e_1(s)=\frac{1}{\sqrt{r^2+a^2}}(-r\sin(s/\sqrt{r^2+a^2}),r\cos(s/\sqrt{r^2+a^2}),a)\\
    e_2(s)=-(\cos(s/\sqrt{r^2+a^2}),\sin(s/\sqrt{r^2+a^2}),0)\\
    e_3(s)=\frac{1}{\sqrt{r^2+a^2}}(a\sin(s/\sqrt{r^2+a^2}),-a\cos(s/\sqrt{r^2+a^2}),r)\\
  \end{cases}
\]
です。

両者で大きく違うのは、\textbf{$e_3$が定値かそうじゃないかです}。
そういえば放物線の$e_3$も
\[
  e_3(s)=(0,0,1)
\]
で定値で、放物線もねじれている感じはありません。

逆にもし$e_3$が定値でなければ、曲線$\gamma$は$e_3$の方向に移動する方向があるということです。
つまり、$d\gamma(t(s))/ds$と$d^2\gamma(t(s))/ds^2$が張る平面である\textbf{接触平面}(定義\ref{definition:tangent_plane}を参照)から離れていく方向の"力"が加わるわけです。
この感覚が人間に「ねじれていない」と思わしめるようです。

\begin{definition}
  $I$を区間、$(I,\gamma)$を弧長パラメータづけ可能な曲線、$s=s(t)$をその弧長パラメータとする。
  さらに
  \[
    \frac{d^2\gamma(t(s))}{ds^2}\neq0\quad({}^\forall s\in s^{-1}(I))
  \]
  を満たすとする。
  このとき、
  \[
    \tau(s):=\left|\frac{de_3(s)}{ds}\right|
  \]
  を、曲線$(I,\gamma)$の\textbf{捩率}という。
\end{definition}



\section{例題：円、放物線、螺旋の曲率}

円や放物線の従法線ベクトル$e_3$は定値だったので、それぞれ
\[
  \tau(s)=\left|\frac{de_3(s)}{ds}\right|=0
\]
となり、ねじれていないと言えます。

一方、常螺旋の場合
\begin{align*}
  &\frac{de_3(s)}{ds}\\
  =&\frac{d}{ds}\frac{1}{\sqrt{r^2+a^2}}(a\sin(s/\sqrt{r^2+a^2}),-a\cos(s/\sqrt{r^2+a^2}),r)\\
  =&\frac{1}{r^2+a^2}
\end{align*}
ゆえに捩率は
\[
  \kappa(s)=\frac{1}{r^2+a^2}
\]
となります。
だいぶねじれてますね。
